% Research questions: each is a numbered, styled item that can be
% cited with \ref (it prints RQ1, RQ2, ...).
\chapter{Research questions}\label{ch:questions}

The gap described above leads to one overarching aim: \emph{to make the
placement of urban heat sensors an explicit, learnable decision that
weights error by human exposure.} Three questions break this aim into
parts that can each be answered within a thesis chapter.

\begin{researchquestion}{Exposure-weighted acquisition}\label{rq:acq}
How should a sequential placement rule trade off reduction in predictive
variance against the number of vulnerable residents near each candidate
site, and how much does such a rule improve mapping error where it
matters most compared with variance-only and random placement?
\end{researchquestion}

\begin{researchquestion}{Transfer between cities}\label{rq:transfer}
Can a temperature model and its placement policy, learned from a dense
network in one city, guide the first deployment in a city with almost
no street-level data, and which covariates make that transfer succeed
or fail?
\end{researchquestion}

\begin{researchquestion}{Field validation}\label{rq:field}
In a live season-long deployment of mobile sensors that are moved every
two weeks, does adaptive placement deliver the gains predicted in
simulation once sensor drift, vandalism and access constraints are
included?
\end{researchquestion}

The questions build on each other: \ref{rq:acq} supplies the rule,
\ref{rq:transfer} asks how far it generalises, and \ref{rq:field} tests
it against the untidiness of a real street. Each question has a
falsifiable outcome, stated as a hypothesis in Chapter~\ref{ch:methods}.
